\documentclass{article}

\begin{document}

\author{V. E. McHale}
\title{Array Language Compilers for SIMD}
\maketitle

\begin{abstract}
  The C compilation model has dominated compilers literature but it opposes Single-Intruction, Multiple Data (SIMD). We articulate the assumptions behind textbook methods, where new CPUs diverge, and give a first cut of a solution with examples from the Apple Array System. By reconsidering expectations about compiler optimization, we give a program for high-level languages to offer performant SIMD. Array languages turn out to be an especially good candidate; we point out that it is possible for the programmer to expect good performance with examples in Apple.
\end{abstract}

% array languages are a good candidate

\section{Background}

\subsection{The C Compilation Model}

Kell \cite{kell2017} points out that C implementations, in practice, offer something that

% elegance: compile code/function to assembly
% https://fabiensanglard.net/dc/

\section{Functions in Theory}

% log is not a "parallel" function b/c conditional... hardware detail?

\section{Compiler Users}

% wrong to expect compiler to optimize code (widespread!), rather, one wants languages/compilers that are amenable to performance.
% really "no big hammers" theorem, not a terrible fundamental... also nothing says high-level can't be performant!
% no "spiritual optimization" but our problem at hand is actually a programmer sitting, writing code. As long as our burden is no more than writing assembly, we offer something as a high-level language
% actually the problem with C++ is more that compilers are fickle, clang vs. gcc https://marek.ai/matrix-multiplication-on-cpu.html

\subsection{Typing}

% indexing by for-loops is elegant but index information (parity!) doesn't propagate so easily (optimization is "internal" rather than specified system)

\bibliographystyle{plain}
\bibliography{simd.bib}

\end{document}
